\documentclass[12pt,a4paper]{article}

\usepackage[spanish]{babel}
\usepackage[utf8]{inputenc}
\usepackage[T1]{fontenc}
\usepackage{lmodern}
\usepackage{setspace}
\usepackage{geometry}
\usepackage{amsmath}
\usepackage{hyperref}

\geometry{margin=2.5cm}
\setstretch{1.2}

\title{Unidad Ontológica Social, Financiarización y Confianza: \\
Hacia una Reconfiguración del Sujeto Histórico de Transformación}
\author{}
\date{}

\begin{document}

\maketitle

\begin{abstract}
Este texto propone una reinterpretación del sujeto histórico de transformación en el capitalismo contemporáneo a partir de tres desplazamientos analíticos fundamentales: la mutación del capitalismo industrial hacia su forma financiarizada, la reconfiguración de la alienación desde su manifestación productiva hacia su interiorización psicosocial, y la comprensión de la sociedad como una unidad ontológica real cuya conciencia no necesita ser producida, sino liberada de los mecanismos que inhiben su expresión. Se argumenta que la propiedad privada funciona hoy como infraestructura operativa del régimen de acumulación financiera y que el obstáculo fundamental para la materialización del sujeto histórico contemporáneo no es la ausencia de conciencia social, sino el miedo que inhibe la confianza. La transformación histórica depende, por tanto, de la institucionalización de condiciones que premien la confianza social.
\end{abstract}

\section{Introducción}

El materialismo histórico clásico identificó en el proletariado industrial el sujeto histórico de transformación del capitalismo del siglo XIX. Este diagnóstico, articulado por Marx en el contexto de la industrialización europea, respondía a una configuración específica de las relaciones de producción, caracterizada por la centralidad de la fábrica, la extracción directa de plusvalía del trabajo asalariado y la visibilidad estructural del antagonismo entre capital y trabajo.

Sin embargo, el capitalismo contemporáneo ha experimentado una mutación profunda. Sus raíces liberalistas permanecen, pero sus mecanismos de acumulación han sido progresivamente desplazados desde la producción material hacia la mediación financiera. Este desplazamiento no elimina la explotación, sino que la reorganiza. La apropiación de valor ya no se limita al tiempo de trabajo presente, sino que se extiende al tiempo vital futuro mediante el endeudamiento estructural.

Este texto parte de la premisa de que dicha transformación exige una reconsideración radical del sujeto histórico. No se trata simplemente de identificar un nuevo grupo social o una nueva clase, sino de comprender la forma en que el capitalismo contemporáneo ha universalizado la interdependencia sistémica hasta hacer coincidir la contradicción estructural con la totalidad del cuerpo social.

\section{De la Producción a la Financiarización}

En el capitalismo industrial analizado por Marx, la acumulación se articulaba fundamentalmente mediante la explotación directa del trabajo vivo. La plusvalía emergía de la diferencia entre el valor producido por el trabajador y el salario recibido.

En el capitalismo financiarizado, la lógica de acumulación se desplaza hacia la extracción de valor mediante el crédito y el interés. La deuda se convierte en el mecanismo central de captura económica. El capital ya no depende exclusivamente de la apropiación del tiempo de trabajo presente, sino que coloniza el tiempo económico futuro de los sujetos.

Este proceso transforma la forma de la subordinación. El trabajador no es únicamente productor de valor, sino también portador de deuda. La explotación deja de estar confinada al espacio productivo y se extiende a la totalidad de la vida social.

\section{La Propiedad Privada como Infraestructura Financiera}

La propiedad privada ha sido históricamente presentada como fundamento moral del liberalismo. Sin embargo, desde una perspectiva materialista, su existencia depende de condiciones contingentes: la capacidad de defensa directa o el reconocimiento institucional sostenido por estructuras jurídicas.

La propiedad no es una entidad ontológicamente estable, sino una relación sostenida por configuraciones históricas de poder. Su carácter aparentemente sagrado no se explica por su estabilidad intrínseca, sino por su función sistémica.

En el capitalismo contemporáneo, la propiedad individualizable constituye la condición de posibilidad del endeudamiento estructural. La deuda requiere sujetos jurídicamente identificables, derechos transferibles y garantías ejecutables. La propiedad privada funciona así como infraestructura operativa del régimen de acumulación financiera.

\section{Transformación de la Alienación}

La alienación descrita por Marx en la producción industrial no desaparece en el capitalismo contemporáneo, sino que se transforma. La subordinación visible del trabajo se convierte en interiorización psicológica del conflicto estructural.

El malestar social se expresa ahora como burnout, ansiedad o crisis existencial. La contradicción sistémica se individualiza. El conflicto deja de percibirse como antagonismo estructural y se experimenta como insuficiencia personal.

Esta interiorización de la alienación constituye una forma particularmente eficaz de estabilización sistémica.

\section{Ontología Social de la Unidad}

La sociedad humana no puede entenderse como mera agregación de individuos autónomos. La interdependencia estructural constituye la condición misma de la existencia humana. La cooperación no es opcional, sino constitutiva.

Esta interdependencia no es una construcción ideológica, sino una realidad material vivida permanentemente. La estabilidad de cualquier proyecto individual depende de la estabilidad del sistema social total.

Incluso fenómenos como la disuasión nuclear global evidencian la conciencia práctica de destino compartido. La humanidad actúa como sistema cuando su supervivencia total está en juego.

\section{El Sujeto Histórico Contemporáneo}

Si el capitalismo industrial produjo un antagonismo localizado entre burguesía y proletariado, el capitalismo financiarizado produce una interdependencia sistémica que integra a toda la sociedad en una estructura común de vulnerabilidad.

El sujeto histórico de transformación ya no puede ser una clase particular. Es la sociedad como totalidad estructural.

Esta totalidad no necesita producir conciencia de sí misma. La conciencia de interdependencia está inscrita biológica y sociológicamente en la existencia humana. El obstáculo no es cognitivo, sino afectivo.

\section{El Miedo como Principio de Fragmentación}

La fragmentación social no deriva de la ignorancia de la unidad, sino del miedo a la vulnerabilidad que implica confiar. La confianza expone, y la exposición genera riesgo.

Las estructuras sociales que mantienen la competencia, la incertidumbre y la responsabilidad individualizada refuerzan este miedo. La desconfianza estabiliza la fragmentación operativa del sistema social.

\section{La Confianza como Principio Transformador}

Si la unidad social es ontológicamente real pero operativamente inhibida, la transformación histórica no depende de producir conciencia colectiva, sino de crear condiciones que reduzcan el coste existencial de confiar.

La acción política fundamental consiste en institucionalizar mecanismos que premien la confianza. La confianza no como valor moral abstracto, sino como principio organizador de la coordinación social.

\section{Dinámica Sistémica de la Confianza y Complejidad Social}

Hasta ahora, la confianza ha sido considerada como principio organizador de la
coordinación social. Sin embargo, para comprender plenamente su función estructural,
es necesario situarla en el marco de la teoría de sistemas complejos.

Desde una perspectiva ontológica ampliada, la confianza puede entenderse como una
relación de acoplamiento correlacional no lineal entre componentes de un sistema.
No se trata únicamente de una disposición psicológica ni de una norma cultural,
sino de una propiedad emergente de sistemas interdependientes cuya estabilidad
depende de la covariación funcional de sus partes.

En este sentido, la sociedad puede interpretarse como una red de acoplamientos
dinámicos entre unidades humanas, institucionales y materiales. La cooperación,
la coordinación y la reproducción social no son simplemente resultados de decisiones
individuales, sino manifestaciones macroscópicas de correlaciones sistémicas
subyacentes.

La confianza no puede ser eliminada estructuralmente en sistemas de este tipo.
Como ocurre con otros fenómenos emergentes en dinámicas no lineales, solo puede
ser modulada localmente mediante intervenciones que alteran las condiciones
de interacción entre los componentes del sistema.

\section{Presión Social y Regulación Sistémica}

Las restricciones estructurales del acoplamiento social generan acumulaciones
de tensión relacional que pueden conceptualizarse como presión social. Esta
presión emerge de la desigualdad persistente, la deuda estructural, la competencia
generalizada, la frustración aspiracional y la limitación de acceso a recursos.

Los sistemas de poder operan mediante la modulación de esta presión acumulada.
Políticas fiscales compensatorias, ciclos de consumo masivo, eventos económicos
ritualizados, narrativas de movilidad futura y dispositivos de regulación institucional
funcionan como mecanismos de liberación controlada de tensión sistémica.

El capitalismo contemporáneo puede entenderse así como un régimen de gestión
dinámica de presión social. Su estabilidad depende de la capacidad de redistribuir
y canalizar continuamente las tensiones generadas por sus propias estructuras
de acumulación.

\section{Control Imperfecto y No Linealidad Social}

Sin embargo, ningún sistema complejo puede ser regulado de forma perfecta.
La dinámica no lineal introduce sensibilidad estructural a perturbaciones,
retroalimentaciones imprevistas y bifurcaciones emergentes.

Las intervenciones diseñadas para estabilizar el sistema pueden producir efectos
globales inesperados. La historia social muestra repetidamente la aparición de
transiciones abruptas que no pueden explicarse únicamente por la intención de
los actores, sino por la dinámica intrínseca del sistema.

Las transformaciones históricas pueden interpretarse, por tanto, como transiciones
de fase emergentes derivadas de la incapacidad del sistema de control para estabilizar
las correlaciones sociales bajo condiciones cambiantes de presión.

\section{Morfogénesis de las Transiciones Históricas}

Cuando el régimen de control pierde eficacia, la reorganización sistémica no es
indeterminada. La forma específica que adopta la nueva fase depende de la capacidad
adaptativa de los remanentes de los estamentos de control preexistentes.

Las élites económicas, institucionales y administrativas intentan reconfigurar sus
estructuras para continuar sus dinámicas fundamentales de acumulación bajo nuevas
condiciones sistémicas. La historia del capitalismo puede entenderse como una
secuencia de reconfiguraciones adaptativas tras episodios de inestabilidad estructural.

La nueva fase histórica no emerge ex nihilo, sino como resultado de la interacción
entre dinámica sistémica emergente y capacidad adaptativa de las estructuras
heredadas.

\section{Confianza, Complejidad y Transformación}

En este marco, la transformación social no consiste en la creación de nuevas
relaciones desde cero, sino en la reorganización global de patrones de acoplamiento
existentes cuando los mecanismos de modulación de presión dejan de ser eficaces.

La confianza, entendida como correlación sistémica efectiva, constituye la base
ontológica de toda reorganización histórica. Su expresión puede ser modulada,
fragmentada o canalizada, pero no suprimida definitivamente.

La evolución social aparece así como un proceso dinámico de modulación, acumulación
de presión, pérdida de control paramétrico y reorganización estructural del sistema
de correlaciones sociales.

\section{Conclusión}

El capitalismo contemporáneo ha universalizado la interdependencia social hasta hacer coincidir la contradicción estructural con la totalidad del cuerpo social. El sujeto histórico de transformación ya no es una clase particular, sino la sociedad como unidad ontológica real.

La conciencia de esta unidad ya existe. Lo que impide su expresión es el miedo que inhibe la confianza. La transformación histórica depende de la creación de estructuras que permitan que la unidad social latente se manifieste como sujeto histórico efectivo.

\section*{Referencias Conceptuales}

Marx, Karl. Crítica de la economía política; El Capital.

Foucault, Michel. Microfísica del poder; Nacimiento de la biopolítica.

Durkheim, Émile. La división del trabajo social.

Polanyi, Karl. La gran transformación.

Luhmann, Niklas. Confianza.

\end{document}
